{\small\opmaakproza
\begin{enumerate}[leftmargin=1.5em,itemsep=3pt,label=\textcolor{kleur1}{\sffamily\bfseries\arabic*}]
\item \textbf{Lees de inleiding van een opgave goed.} Gegevens die je bij een latere vraag nodig hebt, staan vaak
  al bovenaan de opgave, in een tekst, tabel of grafiek.
\item \textbf{Let op het werkwoord.} Bij \emph{bereken} mag je de rekenmachine gebruiken, maar je schrijft op wat je
  doet. Bij \emph{toon aan} werk je zichtbaar naar het gegeven antwoord toe. Bij \emph{beredeneer} of \emph{leg uit}
  telt de redenering, niet alleen het antwoord.
\item \textbf{Schrijf op wat je rekenmachine doet.} Noteer de formule of vergelijking die je invoert en welke optie je
  gebruikt (tabel, snijpunt, normale verdeling met $\mu$, $\sigma$ en de grenzen).
\item \textbf{Procenten en groei.} Bij een toename van $p\%$ hoort de groeifactor $1 + \frac{p}{100}$. Let goed op de
  tijdseenheid: per jaar, per maand of per dag.
\item \textbf{Lees grafieken en tabellen nauwkeurig.} Kijk naar de schaal en de eenheden op de assen, ook bij een
  logaritmische schaal of een grafiek met cumulatieve percentages.
\item \textbf{Gebruik het formuleblad.} Daarop staan de vuistregels voor het verschil tussen twee groepen (phi,
  max $V_{cp}$, effectgrootte, boxplots) en het betrouwbaarheidsinterval. Geef je conclusie in woorden:
  \emph{het verschil is groot}, \emph{middelmatig} of \emph{gering}.
\item \textbf{Rond pas aan het eind af}, en precies zoals de vraag vraagt. Denk aan de context: een aantal mensen is
  een geheel getal.
\item \textbf{Gebruik je eigen antwoord.} Kom je bij een vraag niet uit, werk dan bij de volgende vraag verder met een
  waarde die je aanneemt en schrijf dat erbij.
\item \textbf{Kijk naar de punten.} Een vraag van 5 punten vraagt meer stappen dan een vraag van 3. Houd ongeveer
  2 minuten per punt aan en ga door als je vastloopt.
\item \textbf{Controleer je antwoord}: kan het kloppen in de context, en staat de eenheid erbij?
\end{enumerate}\par}
